\documentclass[../DoAn.tex]{subfiles}
\begin{document}

% Đặt tên caption bảng bằng tiếng Việt khi biên dịch riêng chương này.
\renewcommand{\tablename}{Bảng}

Chương này trình bày phương pháp đề xuất và thiết kế hệ thống mô phỏng của đồ án. Nếu Chương 2 đã xây dựng cơ sở lý thuyết cho bài toán bao phủ trên bản đồ lưới có trọng số, thì Chương 3 tập trung vào cách các khái niệm đó được hiện thực hóa thành một quy trình mô phỏng cụ thể. Nội dung chương bao gồm cấu trúc dữ liệu bản đồ, mô hình trạng thái robot, cách tính chi phí trong hệ thống, thuật toán Dijkstra có hướng, chiến lược chọn mục tiêu bao phủ, cơ chế quay về căn cứ, xử lý vật cản động, vòng lặp mô phỏng và hệ thống ghi nhận kết quả.

\section{Tổng quan phương pháp đề xuất}

Phương pháp đề xuất trong đồ án là xây dựng một hệ thống mô phỏng robot tự hành thực hiện nhiệm vụ bao phủ trên bản đồ lưới có trọng số. Hệ thống xét đồng thời các yếu tố: vật cản tĩnh, vật cản động, chi phí địa hình, hướng hiện tại của robot, năng lượng di chuyển, năng lượng quay, khả năng quay về căn cứ và trạng thái kết thúc nhiệm vụ. Cách thiết kế này hướng tới mục tiêu không chỉ làm cho robot đi qua nhiều ô nhất có thể, mà còn giúp robot duy trì khả năng vận hành an toàn trong quá trình thực hiện nhiệm vụ.

Luồng xử lý tổng quát của hệ thống gồm các bước chính được trình bày trong Bảng~\ref{tab:overall_simulation_flow}.

\begin{table}[!htbp]
\centering
\caption{Quy trình tổng quát của hệ thống mô phỏng}
\label{tab:overall_simulation_flow}
\small
\setlength{\tabcolsep}{4pt}
\renewcommand{\arraystretch}{1.05}
\begin{tabular}{|C{0.10\textwidth}|L{0.26\textwidth}|L{0.54\textwidth}|}
\hline
\textbf{Bước} & \textbf{Pha xử lý} & \textbf{Nội dung thực hiện} \\
\hline
1 & Đọc dữ liệu đầu vào & Hệ thống đọc dung lượng pin, kích thước bản đồ, ma trận địa hình, danh sách vật cản động và tọa độ căn cứ. \\
\hline
2 & Khởi tạo mô phỏng & Khởi tạo robot, căn cứ, tập ô đã bao phủ, vật cản tĩnh, vật cản động và trạng thái nhiệm vụ ban đầu. \\
\hline
3 & Cập nhật môi trường & Ở mỗi bước thời gian, hệ thống cập nhật vị trí vật cản động và các ô đang tạm thời không an toàn. \\
\hline
4 & Lập kế hoạch & Nếu robot chưa có đường đi hiện tại, hệ thống chọn mục tiêu bao phủ và tìm đường bằng Dijkstra có hướng. \\
\hline
5 & Kiểm tra năng lượng & Robot chỉ chọn mục tiêu nếu năng lượng còn lại đủ để đi tới mục tiêu, quay về căn cứ và giữ phần dự phòng khi kịch bản có vật cản động. \\
\hline
6 & Thực thi chuyển động & Nếu ô tiếp theo an toàn, robot thực hiện quay, di chuyển, tiêu hao năng lượng và cập nhật ô đã bao phủ. \\
\hline
7 & Ứng phó khi bị chặn & Nếu ô tiếp theo không an toàn, robot chờ, chuyển sang trạng thái cảnh báo hoặc lập kế hoạch lại. \\
\hline
8 & Ghi nhận kết quả & Hệ thống cập nhật log, hiển thị mô phỏng, thống kê và dữ liệu benchmark. \\
\hline
\end{tabular}
\end{table}

Hệ thống hoạt động theo vòng lặp nhiệm vụ. Ở mỗi bước, robot không chỉ chọn mục tiêu gần nhất, mà còn xét chi phí đường đi, năng lượng còn lại, đường quay về căn cứ và trạng thái an toàn của môi trường. Với kịch bản tĩnh, robot chỉ cần đủ năng lượng để đi tới mục tiêu và quay về căn cứ. Với kịch bản có vật cản động, hệ thống giữ thêm quỹ dự phòng quay về vì đường về có thể phát sinh chi phí do chờ, đi vòng hoặc lập kế hoạch lại. Nhờ đó, chương trình mô phỏng được các ràng buộc chính của một nhiệm vụ dài: phải bao phủ thêm ô mới, nhưng vẫn bảo toàn khả năng quay về và xử lý vật cản động.

Các thành phần chính của hệ thống được tóm tắt trong Bảng~\ref{tab:system_components}.

\begin{table}[!htbp]
\centering
\caption{Các thành phần chính trong hệ thống mô phỏng}
\label{tab:system_components}
\small
\setlength{\tabcolsep}{4pt}
\renewcommand{\arraystretch}{1.05}
\begin{tabular}{|L{0.28\textwidth}|L{0.62\textwidth}|}
\hline
\textbf{Thành phần} & \textbf{Vai trò trong hệ thống} \\
\hline
Bản đồ lưới có trọng số & Biểu diễn môi trường, vật cản tĩnh và chi phí năng lượng tương đối của từng ô hợp lệ. \\
\hline
Robot mô phỏng & Lưu vị trí, hướng, năng lượng, đường đi hiện tại, số bước, số lần quay về căn cứ và các thống kê nhiệm vụ. \\
\hline
Bộ lập kế hoạch & Tìm đường đi chi phí nhỏ nhất bằng Dijkstra có hướng trên trạng thái gồm vị trí và hướng. \\
\hline
Bộ chọn mục tiêu & Chọn ô chưa bao phủ có thể đi tới và vẫn bảo đảm khả năng quay về căn cứ theo ràng buộc năng lượng. \\
\hline
Bộ xử lý năng lượng & Tính năng lượng di chuyển, năng lượng quay, dự phòng quay về có điều kiện và điều kiện quay về căn cứ. \\
\hline
Bộ xử lý vật cản động & Cập nhật các ô tạm thời không an toàn, phát hiện đường bị chặn và kích hoạt chờ hoặc lập kế hoạch lại. \\
\hline
Bộ ghi nhận kết quả & Ghi log, thống kê, dữ liệu benchmark và ảnh trạng thái cuối để phục vụ đánh giá thực nghiệm. \\
\hline
\end{tabular}
\end{table}

\section{Biểu diễn dữ liệu bản đồ đầu vào}

Hệ thống sử dụng bản đồ lưới hai chiều, trong đó mỗi ô có thể là ô hợp lệ hoặc ô vật cản tĩnh. Khác với mô hình lưới chỉ gồm ô trống và ô bị chặn, bản đồ trong đồ án cho phép mỗi ô hợp lệ mang một trọng số địa hình dương. Trọng số này chính là chi phí năng lượng tương đối đã chuẩn hóa như đã trình bày ở Chương 2.

Trong triển khai, hệ thống chỉ yêu cầu các chi phí địa hình giữ đúng quan hệ tương đối giữa các loại ô. Điều này có nghĩa là bản đồ đầu vào không nhất thiết phải dùng trực tiếp đơn vị năng lượng vật lý như Joule. Các giá trị có thể là điểm chi phí đã chuẩn hóa, miễn là ô khó đi có chi phí lớn hơn ô dễ đi theo đúng giả định mô phỏng. Cách biểu diễn này giúp hệ thống dễ tạo bản đồ thử nghiệm, đồng thời vẫn giữ được ý nghĩa tối ưu năng lượng tương đối.

Dữ liệu đầu vào của một bản đồ gồm các thành phần được trình bày trong Bảng~\ref{tab:input_data_structure}.

\begin{table}[!htbp]
\centering
\caption{Cấu trúc dữ liệu đầu vào của hệ thống mô phỏng}
\label{tab:input_data_structure}
\small
\setlength{\tabcolsep}{4pt}
\renewcommand{\arraystretch}{1.05}
\begin{tabular}{|L{0.28\textwidth}|L{0.62\textwidth}|}
\hline
\textbf{Thành phần đầu vào} & \textbf{Ý nghĩa} \\
\hline
Dung lượng pin ban đầu & Giá trị năng lượng tối đa của robot trong một lần sạc. \\
\hline
Số hàng và số cột & Kích thước bản đồ lưới. \\
\hline
Ma trận địa hình & Mỗi ô là một số nguyên dương biểu diễn chi phí địa hình, hoặc ký hiệu vật cản tĩnh. \\
\hline
Danh sách vật cản động & Mỗi vật cản động có loại và tọa độ xuất phát trên bản đồ. \\
\hline
Tọa độ robot xuất phát & Vị trí ban đầu của robot, đồng thời được dùng làm căn cứ. \\
\hline
Tùy chọn mô hình toàn hướng (\texttt{metric\_h}) & Nếu được bật, hệ thống bỏ qua chi phí quay để tạo mô hình đối chứng. \\
\hline
\end{tabular}
\end{table}

Trong quá trình đọc dữ liệu, các ô vật cản tĩnh được gán chi phí vô cùng lớn để loại khỏi quá trình tìm đường. Các ô hợp lệ phải có chi phí dương. Robot và vật cản động không được khởi tạo trên ô vật cản. Ngoài ra, vật cản động cũng không được khởi tạo trùng với căn cứ của robot. Những kiểm tra này giúp bảo đảm bản đồ đầu vào hợp lệ trước khi mô phỏng bắt đầu.

Gọi bản đồ có $m$ hàng và $n$ cột. Mỗi ô được biểu diễn bởi cặp tọa độ theo Công thức~\ref{eq:grid_cell_coordinate}.

\begin{equation}
    v=(r,c),\quad 1 \le r \le m,\quad 1 \le c \le n.
    \label{eq:grid_cell_coordinate}
\end{equation}

Tại mỗi ô, hệ thống lưu các thông tin chính theo Công thức~\ref{eq:cell_state_model}.

\begin{equation}
    \text{cell}(v)=\bigl(c(v),\ \text{staticBlocked}(v),\ \text{dynamicBlocked}_t(v),\ \text{covered}(v)\bigr).
    \label{eq:cell_state_model}
\end{equation}

Trong đó $c(v)$ là chi phí địa hình của ô, $\text{staticBlocked}(v)$ cho biết ô có phải vật cản tĩnh hay không, $\text{dynamicBlocked}_t(v)$ cho biết ô có đang bị vật cản động chiếm tại thời điểm $t$ hay không, còn $\text{covered}(v)$ cho biết robot đã bao phủ ô này hay chưa. Vật cản tĩnh là ràng buộc cố định, còn vật cản động là ràng buộc thay đổi theo thời gian.

Căn cứ của robot được ký hiệu là $b$. Trong hệ thống, căn cứ chính là ô xuất phát của robot, như Công thức~\ref{eq:base_cell_definition}.

\begin{equation}
    b = v_0.
    \label{eq:base_cell_definition}
\end{equation}

Robot có thể quay về căn cứ để sạc pin hoặc kết thúc nhiệm vụ. Việc xem ô xuất phát là căn cứ giúp mô hình giữ được tính đơn giản nhưng vẫn đủ để khảo sát cơ chế quay về và phục hồi năng lượng.

Khi khởi tạo robot, hệ thống không cố định thủ công hướng ban đầu. Thay vào đó, bốn hướng khả dĩ được đánh giá trước nhiệm vụ. Với mỗi hướng, hệ thống chạy Dijkstra có hướng trong chế độ bỏ qua vật cản động để xác định số mục tiêu bao phủ có thể tiếp cận, chi phí trung bình tới các mục tiêu đó và số ô trống liên tiếp theo hướng thẳng ban đầu. Hướng ban đầu được chọn theo thứ tự ưu tiên: tiếp cận được nhiều mục tiêu hơn, chi phí trung bình thấp hơn, đi thẳng được xa hơn, sau đó mới dùng thứ tự hướng làm tiêu chí phụ. Cách chọn này tránh việc đưa vào một giới hạn mẫu tùy ý, đồng thời vẫn chỉ làm tăng chi phí khởi tạo thêm bốn lần chạy Dijkstra có hướng.

Trong phiên bản hiện tại, hệ thống sử dụng một căn cứ duy nhất để tập trung vào cơ chế quay về, sạc và tiếp tục nhiệm vụ. Nếu mở rộng sang nhiều trạm sạc hoặc nhiều điểm thu hồi, có thể thay căn cứ đơn $b$ bằng một tập căn cứ $B$ và chọn điểm quay về có chi phí khả thi nhỏ nhất. Phần mở rộng đó không làm thay đổi bản chất của mô hình lập kế hoạch, nhưng không được triển khai trong phạm vi hiện tại để giữ hệ thống gọn và dễ đánh giá.

\section{Mô hình trạng thái robot}

Robot trong hệ thống được mô phỏng ở mức lập kế hoạch và điều khiển rời rạc trên bản đồ lưới. Tại thời điểm $t$, trạng thái chính của robot có thể được biểu diễn bằng Công thức~\ref{eq:robot_state_model}.

\begin{equation}
    R_t = (v_t,\ h_t,\ E_t,\ P_t,\ mode_t).
    \label{eq:robot_state_model}
\end{equation}

Trong đó $v_t$ là vị trí hiện tại, $h_t$ là hướng hiện tại, $E_t$ là năng lượng còn lại, $P_t$ là đường đi hiện tại mà robot đang theo, và $mode_t$ là trạng thái nhiệm vụ của robot. Ngoài các thành phần trên, hệ thống còn lưu số bước đã đi, tổng năng lượng đã dùng, năng lượng dùng cho di chuyển, năng lượng dùng cho quay, số lần quay về căn cứ, số lần sạc, các cạnh đã đi qua và vết đường di chuyển của robot.

Tập hướng của robot gồm bốn hướng chính, được định nghĩa trong Công thức~\ref{eq:direction_set}.

\begin{equation}
    H=\{N,E,S,W\}.
    \label{eq:direction_set}
\end{equation}

Do đó, một trạng thái dùng cho lập kế hoạch đường đi có dạng như Công thức~\ref{eq:oriented_state}.

\begin{equation}
    q_t=(v_t,h_t).
    \label{eq:oriented_state}
\end{equation}

Việc lưu hướng hiện tại là điểm khác biệt quan trọng so với mô hình chỉ xét tọa độ. Nếu hai trạng thái có cùng vị trí nhưng khác hướng, chi phí đi tiếp sang cùng một ô kề cạnh có thể khác nhau do số lần quay khác nhau. Đây là cơ sở để hệ thống sử dụng Dijkstra có hướng thay vì Dijkstra chỉ trên tọa độ ô.

Trạng thái nhiệm vụ của robot được mô tả bằng các mode chính trong Bảng~\ref{tab:robot_modes}.

\begin{table}[!htbp]
\centering
\caption{Các trạng thái nhiệm vụ của robot trong hệ thống}
\label{tab:robot_modes}
\small
\setlength{\tabcolsep}{4pt}
\renewcommand{\arraystretch}{1.05}
\begin{tabular}{|L{0.3\textwidth}|L{0.6\textwidth}|}
\hline
\textbf{Trạng thái} & \textbf{Ý nghĩa} \\
\hline
NORMAL & Robot đang thực hiện bao phủ bình thường. \\
\hline
ALERT & Robot phát hiện vật cản động gần hoặc đường hiện tại bị ảnh hưởng, cần thận trọng và có thể lập kế hoạch lại. \\
\hline
HOLD\_SAFE & Robot tạm dừng tại vị trí an toàn khi chưa tìm được đường đi phù hợp trong ngắn hạn. \\
\hline
RETURN\_TO\_BASE & Robot đang quay về căn cứ do năng lượng thấp, nhiệm vụ đã bao phủ xong hoặc cần phục hồi trạng thái. \\
\hline
RECHARGING & Robot đã về căn cứ và đang sạc lại năng lượng. \\
\hline
POWER\_SAVE & Robot dừng để bảo toàn trạng thái khi không thể tiếp tục nhiệm vụ một cách an toàn. \\
\hline
WAIT\_FOR\_COMMAND & Robot ở tình huống tới hạn và chờ chỉ thị xử lý tiếp theo trong mô phỏng. \\
\hline
FINAL\_PUSH & Robot tiếp tục nhiệm vụ trong tình huống đặc biệt khi hệ thống được cấp chỉ thị tương ứng. \\
\hline
\end{tabular}
\end{table}

Các trạng thái này giúp mô phỏng tách rõ hành vi bao phủ, hành vi cảnh báo, hành vi quay về căn cứ và hành vi kết thúc nhiệm vụ. Đồ án không giả định robot luôn hoàn thành nhiệm vụ trong mọi trường hợp. Thay vào đó, hệ thống phân biệt các mức kết quả như hoàn thành đầy đủ, hoàn thành một phần nhưng quay về được, bảo toàn được trạng thái nhiệm vụ hoặc thất bại.

\section{Cách tính chi phí trong hệ thống}

Chương 2 đã trình bày mô hình năng lượng tổng quát. Trong hệ thống mô phỏng, mô hình này được cụ thể hóa theo các quy tắc rời rạc để phù hợp với bản đồ lưới bốn hướng.

Khi robot đi từ ô $v$ sang ô kề cạnh $v'$, năng lượng di chuyển chuẩn hóa được tính bằng chi phí đi vào ô $v'$, như Công thức~\ref{eq:move_energy}.

\begin{equation}
    E_{\text{move}}(v,v') = c_{\text{eff}}(v').
    \label{eq:move_energy}
\end{equation}

Ở đây $c_{\text{eff}}(v')$ là chi phí hiệu dụng của ô $v'$ tại thời điểm lập kế hoạch. Nếu ô $v'$ là vật cản tĩnh hoặc đang bị vật cản động chiếm trong chế độ xét vật cản động, ta có Công thức~\ref{eq:blocked_cell_cost}.

\begin{equation}
    c_{\text{eff}}(v') = \infty.
    \label{eq:blocked_cell_cost}
\end{equation}

Ngược lại, nếu ô $v'$ hợp lệ và an toàn tại thời điểm xét, ta có Công thức~\ref{eq:valid_cell_cost}.

\begin{equation}
    c_{\text{eff}}(v') = c(v').
    \label{eq:valid_cell_cost}
\end{equation}

Chi phí quay được tính theo số lần quay 90 độ cần thiết để chuyển từ hướng hiện tại $h$ sang hướng tiếp theo $h'$. Gọi $\Delta(h,h')$ là số lần quay 90 độ nhỏ nhất giữa hai hướng, như Công thức~\ref{eq:turn_delta}.

\begin{equation}
    \Delta(h,h') \in \{0,1,2\}.
    \label{eq:turn_delta}
\end{equation}

Trong mô hình mặc định, chi phí cho một lần quay 90 độ tại ô $v$ được tính theo Công thức~\ref{eq:quarter_turn_cost}.

\begin{equation}
    \tau(v)=\frac{1}{2}c(v).
    \label{eq:quarter_turn_cost}
\end{equation}

Do đó, chi phí chuyển từ trạng thái $(v,h)$ sang trạng thái $(v',h')$ là Công thức~\ref{eq:transition_cost}.

\begin{equation}
    w_q\bigl((v,h),(v',h')\bigr)=c_{\text{eff}}(v')+\Delta(h,h')\tau(v).
    \label{eq:transition_cost}
\end{equation}

Công thức trên là công thức chi phí trực tiếp được dùng trong Dijkstra có hướng. Nó kết hợp chi phí đi vào ô kế tiếp và chi phí quay tại ô hiện tại. Nếu robot đi thẳng theo hướng đang có, $\Delta(h,h')=0$ và chi phí chỉ còn là chi phí di chuyển. Nếu robot cần quay 90 độ, chi phí tăng thêm $\frac{1}{2}c(v)$. Nếu robot quay 180 độ, chi phí tăng thêm $c(v)$.

Trong cấu hình mô hình toàn hướng, chi phí quay được bỏ qua như Công thức~\ref{eq:holonomic_turn_cost}.

\begin{equation}
    \tau(v)=0.
    \label{eq:holonomic_turn_cost}
\end{equation}

Khi đó chi phí chuyển trạng thái rút gọn thành Công thức~\ref{eq:holonomic_transition_cost}.

\begin{equation}
    w_q\bigl((v,h),(v',h')\bigr)=c_{\text{eff}}(v').
    \label{eq:holonomic_transition_cost}
\end{equation}

Bảng~\ref{tab:transition_cost_rules} tóm tắt các quy tắc chi phí được sử dụng trong mô phỏng.

\begin{table}[!htbp]
\centering
\caption{Quy tắc tính chi phí chuyển trạng thái trong hệ thống}
\label{tab:transition_cost_rules}
\small
\setlength{\tabcolsep}{4pt}
\renewcommand{\arraystretch}{1.05}
\begin{tabular}{|L{0.28\textwidth}|L{0.30\textwidth}|L{0.32\textwidth}|}
\hline
\textbf{Thao tác} & \textbf{Mô hình mặc định} & \textbf{Mô hình toàn hướng} \\
\hline
Đi thẳng sang ô $v'$ & $c_{\text{eff}}(v')$ & $c_{\text{eff}}(v')$ \\
\hline
Quay 90 độ rồi đi & $\frac{1}{2}c(v)+c_{\text{eff}}(v')$ & $c_{\text{eff}}(v')$ \\
\hline
Quay 180 độ rồi đi & $c(v)+c_{\text{eff}}(v')$ & $c_{\text{eff}}(v')$ \\
\hline
Ô không an toàn hoặc vật cản & $\infty$ & $\infty$ \\
\hline
\end{tabular}
\end{table}

Các giá trị năng lượng được lượng tử hóa theo bước $0.5$ đơn vị năng lượng. Điều này phù hợp với cách chuẩn hóa chi phí quay 90 độ bằng một nửa chi phí địa hình tại ô hiện tại. Nhờ đó, mô hình có thể biểu diễn được các bước quay 90 độ, quay 180 độ và di chuyển sang ô kế tiếp trên cùng một thang đo chi phí.

\section{Dijkstra có hướng}

Bộ lập kế hoạch của hệ thống sử dụng Dijkstra có hướng. Thay vì tìm đường trên tập ô $V_{\text{free}}$, thuật toán tìm đường trên không gian trạng thái mở rộng được định nghĩa trong Công thức~\ref{eq:directed_state_space}.

\begin{equation}
    Q = V_{\text{free}} \times H.
    \label{eq:directed_state_space}
\end{equation}

Mỗi trạng thái $q=(v,h)$ gồm vị trí $v$ và hướng hiện tại $h$. Từ một trạng thái, thuật toán xét bốn hướng di chuyển có thể có. Với mỗi hướng, hệ thống xác định ô kề cạnh tương ứng, kiểm tra ô đó có hợp lệ không, sau đó tính chi phí chuyển trạng thái theo công thức đã nêu ở mục trước.

Khoảng cách tối ưu tới một trạng thái được lưu dưới dạng Công thức~\ref{eq:state_distance}.

\begin{equation}
    d(r,c,h).
    \label{eq:state_distance}
\end{equation}

Trạng thái bắt đầu có chi phí bằng 0, như Công thức~\ref{eq:start_distance}.

\begin{equation}
    d(r_0,c_0,h_0)=0.
    \label{eq:start_distance}
\end{equation}

Khi xét một bước chuyển từ trạng thái $q$ sang trạng thái $q'$, thuật toán cập nhật theo Công thức~\ref{eq:dijkstra_relaxation}.

\begin{equation}
    d(q') \leftarrow \min\bigl(d(q'),\ d(q)+w_q(q,q')\bigr).
    \label{eq:dijkstra_relaxation}
\end{equation}

Pseudo-code của Dijkstra có hướng trong hệ thống có thể mô tả bằng Bảng~\ref{tab:directed_dijkstra_process}.

\begin{table}[!htbp]
\centering
\caption{Quy trình Dijkstra có hướng trong hệ thống}
\label{tab:directed_dijkstra_process}
\small
\setlength{\tabcolsep}{4pt}
\renewcommand{\arraystretch}{1.05}
\begin{tabular}{|C{0.10\textwidth}|L{0.30\textwidth}|L{0.50\textwidth}|}
\hline
\textbf{Bước} & \textbf{Thao tác} & \textbf{Ý nghĩa} \\
\hline
1 & Khởi tạo bảng khoảng cách & Gán $d(r,c,h)=\infty$ cho mọi trạng thái có hướng và gán chi phí trạng thái bắt đầu bằng 0. \\
\hline
2 & Đưa trạng thái bắt đầu vào hàng đợi ưu tiên & Bắt đầu quá trình tìm đường từ vị trí và hướng hiện tại của robot. \\
\hline
3 & Lấy trạng thái có chi phí nhỏ nhất & Bảo đảm nguyên lý của Dijkstra: trạng thái được xử lý theo thứ tự chi phí không giảm. \\
\hline
4 & Xét bốn hướng di chuyển & Sinh các trạng thái kế tiếp tương ứng với bốn ô kề cạnh trên bản đồ lưới. \\
\hline
5 & Kiểm tra tính hợp lệ của ô kế tiếp & Bỏ qua ô ngoài bản đồ, ô vật cản tĩnh hoặc ô đang không an toàn theo chế độ xét vật cản động. \\
\hline
6 & Tính chi phí chuyển trạng thái & Kết hợp chi phí đi vào ô kế tiếp $c_{\text{eff}}(v')$ và chi phí quay $\Delta(h,h')\tau(v)$. \\
\hline
7 & Cập nhật trạng thái nếu tốt hơn & Nếu chi phí mới nhỏ hơn chi phí đã biết, hệ thống cập nhật bảng khoảng cách và lưu vết truy ngược. \\
\hline
8 & Truy vết đường đi & Sau khi chọn được mục tiêu, hệ thống khôi phục đường đi từ bảng truy vết để robot thực thi. \\
\hline
\end{tabular}
\end{table}

Hệ thống lưu thêm bảng truy vết để khôi phục đường đi sau khi tìm được mục tiêu. Khi cần đường đi tới một ô đích, hệ thống không cố định trước hướng đến đích. Thay vào đó, nó xét bốn hướng có thể tại ô đích và chọn hướng có chi phí nhỏ nhất, như Công thức~\ref{eq:best_goal_orientation}.

\begin{equation}
    d^*(g)=\min_{h\in H} d(g,h).
    \label{eq:best_goal_orientation}
\end{equation}

Cách làm này phù hợp với bài toán bao phủ, vì mục tiêu của robot là đi tới ô cần bao phủ với chi phí thấp nhất, không bắt buộc phải đến ô đó theo một hướng cụ thể. Tuy nhiên, hướng đến đích vẫn được lưu lại vì nó ảnh hưởng tới bước lập kế hoạch tiếp theo, đặc biệt là khi cần ước lượng chi phí từ mục tiêu quay về căn cứ.

\section{Chiến lược chọn mục tiêu bao phủ}

Tại mỗi thời điểm cần lập kế hoạch mới, hệ thống xác định tập các ô chưa bao phủ theo Công thức~\ref{eq:uncovered_set}.

\begin{equation}
    U_t=\{v\in V_{\text{free}} \mid \text{covered}_t(v)=false\}.
    \label{eq:uncovered_set}
\end{equation}

Không phải mọi ô trong $U_t$ đều có thể được chọn làm mục tiêu ngay lập tức. Một ô chỉ được xem là ứng viên nếu robot có đường đi tới ô đó theo trạng thái hiện tại và đường đi này có chi phí hữu hạn. Với mỗi mục tiêu ứng viên $g$, hệ thống tính hai đại lượng trong Công thức~\ref{eq:cost_to_goal} và Công thức~\ref{eq:cost_return_from_goal}.

\begin{equation}
    d_{\text{to}}(g)=E(P_{v_t\rightarrow g}).
    \label{eq:cost_to_goal}
\end{equation}

\begin{equation}
    d_{\text{return}}(g)=E(P_{g\rightarrow b}).
    \label{eq:cost_return_from_goal}
\end{equation}

Trong đó $d_{\text{to}}(g)$ là chi phí từ vị trí hiện tại của robot tới mục tiêu $g$, còn $d_{\text{return}}(g)$ là chi phí ước lượng từ mục tiêu $g$ quay về căn cứ. Cả hai chi phí này đều được tính bằng Dijkstra có hướng.

Hệ thống sắp xếp các ứng viên theo chi phí đi tới mục tiêu. Sau đó, lần lượt kiểm tra từng ứng viên theo điều kiện năng lượng. Gọi $M(d)$ là năng lượng dự phòng cần giữ lại cho đường quay về có chi phí $d$. Mục tiêu $g$ được xem là khả thi với năng lượng hiện tại nếu thỏa Công thức~\ref{eq:target_energy_feasibility}.

\begin{equation}
    E_t \ge d_{\text{to}}(g)+d_{\text{return}}(g)+M(d_{\text{return}}(g)).
    \label{eq:target_energy_feasibility}
\end{equation}

Trong phiên bản hiện tại, $M(d)$ được áp dụng theo loại kịch bản như Công thức~\ref{eq:return_margin_policy}.

\begin{equation}
M(d)=
\begin{cases}
0, & \text{khi kịch bản không có vật cản động},\\
\max\left(M_{\min},M_{\text{prop}}(d)\right), & \text{khi kịch bản có vật cản động}.
\end{cases}
\label{eq:return_margin_policy}
\end{equation}

Cách tách này phản ánh đúng vai trò của quỹ dự phòng. Với bản đồ tĩnh, đường quay về được đánh giá trên bản đồ đã biết nên robot chỉ cần đủ năng lượng cho chặng đi tới mục tiêu và chặng quay về căn cứ. Với kịch bản có vật cản động, hệ thống giữ thêm dự phòng vì đường về có thể phát sinh chi phí do chờ, đi vòng hoặc lập kế hoạch lại.

Pseudo-code của bước chọn mục tiêu có thể viết như Bảng~\ref{tab:target_selection_process}.

\begin{table}[!htbp]
\centering
\caption{Quy trình lựa chọn mục tiêu bao phủ}
\label{tab:target_selection_process}
\small
\setlength{\tabcolsep}{4pt}
\renewcommand{\arraystretch}{1.05}
\begin{tabular}{|C{0.10\textwidth}|L{0.34\textwidth}|L{0.46\textwidth}|}
\hline
\textbf{Bước} & \textbf{Điều kiện kiểm tra} & \textbf{Hành động của hệ thống} \\
\hline
1 & Tồn tại ô chưa bao phủ & Tạo danh sách các ô chưa bao phủ làm ứng viên mục tiêu. \\
\hline
2 & Có đường từ robot tới mục tiêu & Tính chi phí $d_{\text{to}}(g)$ bằng Dijkstra có hướng. \\
\hline
3 & Có đường từ mục tiêu về căn cứ & Tính chi phí $d_{\text{return}}(g)$ bằng Dijkstra có hướng từ mục tiêu về căn cứ. \\
\hline
4 & Pin tối đa không đủ để đi và quay về & Đánh dấu mục tiêu không khả thi dài hạn và xét ứng viên tiếp theo. \\
\hline
5 & Pin hiện tại không đủ để đi và quay về & Đánh dấu nhu cầu quay về sạc trước khi tiếp tục bao phủ. \\
\hline
6 & Đủ năng lượng và đường hiện tại an toàn & Chọn mục tiêu, dựng đường đi và giao cho robot thực thi. \\
\hline
7 & Không có mục tiêu phù hợp tại thời điểm hiện tại & Robot chờ, lập kế hoạch lại, chuyển sang HOLD\_SAFE hoặc quay về căn cứ tùy trạng thái. \\
\hline
\end{tabular}
\end{table}

Điểm quan trọng của chiến lược này là robot không chọn mục tiêu chỉ vì mục tiêu đó gần nhất theo số bước. Mục tiêu được chọn phải đồng thời thỏa mãn ba điều kiện: có đường đi tới, có đường quay về và đủ năng lượng theo chính sách dự phòng của kịch bản. Cách chọn này giúp hạn chế tình huống robot đạt thêm một phần bao phủ trước mắt nhưng làm suy giảm khả năng quay về căn cứ.

Chiến lược trên thuộc nhóm chọn mục tiêu tham lam có ràng buộc. Hệ thống không tìm toàn bộ thứ tự bao phủ tối ưu cho cả nhiệm vụ, mà chọn mục tiêu khả thi tốt nhất tại thời điểm lập kế hoạch hiện tại. Cách làm này giữ quá trình mô phỏng đơn giản, kiểm tra được bằng log và đủ để đánh giá tác động của chi phí quay, trọng số địa hình, năng lượng và vật cản động.

\section{Cơ chế quay về căn cứ và sạc}

Cơ chế quay về căn cứ là một thành phần quan trọng của hệ thống. Robot có thể chuyển sang trạng thái quay về căn cứ trong một số trường hợp: năng lượng còn lại thấp, không còn đủ năng lượng để chọn mục tiêu mới, nhiệm vụ đã bao phủ xong nhưng robot chưa ở căn cứ, hoặc môi trường bị chặn trong thời gian dài và robot cần phục hồi trạng thái.

Tại vị trí hiện tại $v_t$, hệ thống ước lượng chi phí quay về căn cứ theo Công thức~\ref{eq:return_to_base_cost}.

\begin{equation}
    d_b = E(P_{v_t\rightarrow b}).
    \label{eq:return_to_base_cost}
\end{equation}

Robot được yêu cầu quay về nếu năng lượng còn lại nhỏ hơn hoặc bằng chi phí quay về cộng với năng lượng dự phòng của kịch bản, như Công thức~\ref{eq:should_return_condition}.

\begin{equation}
    E_t \le d_b + M(d_b).
    \label{eq:should_return_condition}
\end{equation}

Trong hệ thống, hàm dự phòng được tính theo Công thức~\ref{eq:return_margin_base_policy}.

\begin{equation}
M(d_b)=
\begin{cases}
0, & \text{khi kịch bản không có vật cản động},\\
\max\left(M_{\min},\ M_{\text{prop}}(d_b)\right), & \text{khi kịch bản có vật cản động}.
\end{cases}
\label{eq:return_margin_base_policy}
\end{equation}

Trong nhánh có vật cản động, $M_{\min}$ là dự phòng tối thiểu và $M_{\text{prop}}(d_b)$ là thành phần dự phòng tăng theo chi phí quay về. Với cấu hình hiện tại, $M_{\min}=15$ và thành phần tỷ lệ được lấy xấp xỉ bằng $25\%$ chi phí quay về, tức Công thức~\ref{eq:proportional_return_margin}.

\begin{equation}
    M_{\text{prop}}(d_b)\approx \frac{d_b}{4}.
    \label{eq:proportional_return_margin}
\end{equation}

Giá trị này được lượng tử hóa theo bước $0.5$ để thống nhất với cách hệ thống ghi nhận năng lượng.

Công thức dự phòng này là một chính sách cụ thể của hệ thống, không phải lựa chọn duy nhất có thể có. Điểm chính của phiên bản hiện tại là dự phòng không được áp dụng máy móc cho mọi bản đồ. Nếu kịch bản không có vật cản động, việc giữ thêm $25\%$ năng lượng cho một rủi ro không tồn tại sẽ làm chính sách trở nên quá bảo thủ. Nếu kịch bản có vật cản động, quỹ dự phòng giúp robot còn dư địa năng lượng khi đường về phát sinh chi phí do chờ, đi vòng hoặc lập kế hoạch lại.

Robot rơi vào vùng năng lượng tới hạn nếu thỏa một trong hai điều kiện ở Công thức~\ref{eq:critical_return_energy} hoặc Công thức~\ref{eq:min_energy_condition}.

\begin{equation}
    E_t \le d_b.
    \label{eq:critical_return_energy}
\end{equation}

\begin{equation}
    E_t \le E_{\min}.
    \label{eq:min_energy_condition}
\end{equation}

Trong cấu hình hiện tại, $E_{\min}=3$. Khi rơi vào vùng năng lượng tới hạn, robot không tiếp tục bao phủ như bình thường mà chuyển sang các cơ chế ứng phó thận trọng hơn như chờ, tìm đường vòng, quay về nếu có thể, hoặc bảo toàn trạng thái nhiệm vụ.

Khi robot về tới căn cứ, nếu nhiệm vụ chưa hoàn thành, robot chuyển sang trạng thái sạc. Sau khi sạc đầy, robot quay lại trạng thái bao phủ bình thường và tiếp tục xử lý các ô chưa bao phủ. Nếu toàn bộ bản đồ đã được bao phủ và robot đã ở căn cứ, nhiệm vụ được xem là hoàn thành đầy đủ. Trường hợp robot vừa về tới căn cứ với năng lượng còn lại bằng 0 vẫn được xem là quay về thành công, vì robot đã tới điểm sạc hoặc điểm kết thúc nhiệm vụ.

\section{Mô hình vật cản động và xử lý an toàn}

Trong hệ thống, vật cản động được mô hình hóa như các tác nhân cùng tồn tại trong môi trường. Chúng không được thiết kế như kẻ địch, không truy đuổi robot và không cố tình chặn robot. Vai trò của chúng là làm cho một số ô trên bản đồ tạm thời trở nên không an toàn, từ đó buộc robot phải có khả năng chờ, tránh hoặc lập kế hoạch lại.

Tại thời điểm $t$, gọi $O_t$ là tập ô đang bị vật cản động chiếm dụng. Khi đó, tập ô an toàn được xác định theo Công thức~\ref{eq:safe_cell_set}.

\begin{equation}
    S_t=V_{\text{free}}\setminus O_t.
    \label{eq:safe_cell_set}
\end{equation}

Một bước đi tới ô $v'$ chỉ được phép bắt đầu nếu thỏa Công thức~\ref{eq:safe_next_cell_condition}.

\begin{equation}
    v'\in S_t.
    \label{eq:safe_next_cell_condition}
\end{equation}

Ngoài ra, hệ thống còn kiểm tra lại ô đích trước khi hoàn tất bước di chuyển. Điều này cần thiết vì vật cản động có thể thay đổi trạng thái trong khi robot đang thực hiện thao tác quay hoặc đang đi tới ô tiếp theo. Nếu ô đích trở nên không an toàn trước khi robot hoàn tất bước đi, hệ thống hủy thao tác di chuyển và chuyển sang cơ chế xử lý đường bị chặn.

Vật cản động trong hệ thống gồm hai nhóm chính, được trình bày trong Bảng~\ref{tab:dynamic_obstacle_types}.

\begin{table}[!htbp]
\centering
\caption{Hai nhóm vật cản động trong hệ thống mô phỏng}
\label{tab:dynamic_obstacle_types}
\small
\setlength{\tabcolsep}{4pt}
\renewcommand{\arraystretch}{1.05}
\begin{tabular}{|L{0.24\textwidth}|L{0.66\textwidth}|}
\hline
\textbf{Loại vật cản động} & \textbf{Mô tả hành vi} \\
\hline
Guard & Tác nhân tuần tra cục bộ quanh vị trí tuần tra ban đầu trong một bán kính giới hạn, có các pha chờ, đi ra và quay về. \\
\hline
Vehicle & Tác nhân di chuyển theo hướng tuần tra, có thể chọn nhánh phụ, quay đầu trong hành lang hoặc chờ ngắn khi bị chặn. \\
\hline
\end{tabular}
\end{table}

Lớp cập nhật vật cản động sử dụng bản đồ chiếm dụng để đánh dấu các ô đang bị chặn. Khi vật cản động rời khỏi một ô hoặc đi vào một ô mới, hệ thống đánh dấu các ô bị thay đổi là các ô bẩn. Ở vòng lặp robot, nếu một ô bẩn ảnh hưởng đến đường đi hiện tại, đường đi đó được xóa và robot lập kế hoạch lại. Cách xử lý này giúp tránh việc robot tiếp tục đi theo một đường đã lỗi thời trong môi trường thay đổi.

Khi đường hiện tại bị chặn, robot không lập tức kết luận nhiệm vụ thất bại. Hệ thống áp dụng các mức ứng phó theo thứ tự: chờ trong thời gian ngắn, chuyển sang trạng thái cảnh báo, thử lập kế hoạch lại, tạm dừng an toàn nếu bị chặn nhiều lần, tìm đường vòng khi quay về căn cứ, hoặc quay về căn cứ để phục hồi nếu chờ quá lâu. Cách xử lý này phù hợp với mục tiêu của đồ án: robot cần đối phó được với tình huống bất lợi thay vì chỉ hoạt động trong môi trường tĩnh hoàn toàn.

\section{Vòng lặp mô phỏng và chuyển trạng thái}

Mô phỏng được tổ chức theo các bước thời gian rời rạc. Ở mỗi bước, hệ thống cập nhật trạng thái robot, trạng thái vật cản động, đường đi hiện tại, năng lượng và dữ liệu hiển thị. Vòng lặp chính có thể được mô tả như Bảng~\ref{tab:simulation_loop_priority}.

\begin{table}[!htbp]
\centering
\caption{Thứ tự xử lý trong vòng lặp mô phỏng}
\label{tab:simulation_loop_priority}
\small
\setlength{\tabcolsep}{4pt}
\renewcommand{\arraystretch}{1.05}
\begin{tabular}{|C{0.12\textwidth}|L{0.32\textwidth}|L{0.46\textwidth}|}
\hline
\textbf{Ưu tiên} & \textbf{Trạng thái hoặc điều kiện} & \textbf{Xử lý trong vòng lặp mô phỏng} \\
\hline
1 & Đang có thao tác quay hoặc di chuyển chưa hoàn tất & Tiếp tục pha quay hoặc pha di chuyển, cập nhật năng lượng và vị trí nếu thao tác hoàn tất. \\
\hline
2 & Robot đang ở trạng thái RECHARGING & Xử lý sạc pin, sau đó chuyển về NORMAL nếu nhiệm vụ còn tiếp tục. \\
\hline
3 & Robot ở POWER\_SAVE hoặc WAIT\_FOR\_COMMAND & Xử lý trạng thái bảo toàn hoặc chờ chỉ thị theo chính sách nhiệm vụ. \\
\hline
4 & Đường hiện tại bị vật cản động ảnh hưởng & Xóa đường đi cũ, chuyển trạng thái cảnh báo và chuẩn bị lập kế hoạch lại. \\
\hline
5 & Robot đang RETURN\_TO\_BASE & Ưu tiên xử lý đường về căn cứ, tìm đường vòng hoặc chờ nếu đường về bị chặn. \\
\hline
6 & Robot đang HOLD\_SAFE & Tiếp tục chờ an toàn, thử phục hồi định kỳ hoặc quay về căn cứ nếu chờ quá lâu. \\
\hline
7 & Năng lượng còn lại yêu cầu quay về & Chuyển sang RETURN\_TO\_BASE trước khi tiếp tục chọn mục tiêu bao phủ mới. \\
\hline
8 & Không có đường đi hiện tại & Chọn mục tiêu bao phủ và lập kế hoạch đường đi mới. \\
\hline
9 & Có đường đi và ô tiếp theo an toàn & Thực hiện bước quay, di chuyển, tiêu hao năng lượng và cập nhật bao phủ. \\
\hline
10 & Ô tiếp theo không an toàn & Chờ, chuyển ALERT, HOLD\_SAFE hoặc lập kế hoạch lại tùy mức độ bị chặn. \\
\hline
\end{tabular}
\end{table}

Thiết kế này tách riêng hai loại hành động: hành động lập kế hoạch và hành động thực thi chuyển động. Khi robot đã bắt đầu một bước di chuyển, hệ thống theo dõi pha quay và pha đi tới ô tiếp theo. Năng lượng quay được tiêu thụ trong pha quay, còn năng lượng di chuyển được tiêu thụ khi robot hoàn tất bước đi vào ô mới. Cách làm này giúp mô phỏng thể hiện rõ sự khác biệt giữa chi phí quay và chi phí di chuyển, đồng thời giúp cấu hình toàn hướng, tức Metric H trong log và mã kịch bản, có hành vi khác biệt: trong cấu hình này, robot không mất năng lượng quay và không mất thời gian mô phỏng cho pha quay.

Các chuyển trạng thái của robot có thể tóm tắt như sau:

\begin{equation}
    NORMAL \rightarrow ALERT
    \label{eq:mode_normal_to_alert}
\end{equation}

khi có vật cản động gần hoặc đường hiện tại bị chặn;

\begin{equation}
    NORMAL \rightarrow RETURN\_TO\_BASE
    \label{eq:mode_normal_to_return}
\end{equation}

khi năng lượng còn lại không đủ an toàn để tiếp tục bao phủ;

\begin{equation}
    ALERT \rightarrow HOLD\_SAFE
    \label{eq:mode_alert_to_hold_safe}
\end{equation}

khi robot nhiều lần không tìm được bước đi an toàn;

\begin{equation}
    RETURN\_TO\_BASE \rightarrow RECHARGING
    \label{eq:mode_return_to_recharging}
\end{equation}

khi robot đã về căn cứ nhưng nhiệm vụ chưa hoàn tất;

\begin{equation}
    RECHARGING \rightarrow NORMAL
    \label{eq:mode_recharging_to_normal}
\end{equation}

khi robot đã sạc đầy;

\begin{equation}
    RETURN\_TO\_BASE \rightarrow MISSION\_SUCCESS
    \label{eq:mode_return_to_success}
\end{equation}

khi robot đã bao phủ xong và quay về căn cứ;

\begin{equation}
    HOLD\_SAFE \rightarrow RETURN\_TO\_BASE
    \label{eq:mode_hold_safe_to_return}
\end{equation}

khi chờ quá lâu mà chưa khôi phục được đường bao phủ phù hợp.

Những chuyển trạng thái này là chính sách nhiệm vụ của chương trình mô phỏng. Chúng không đóng vai trò bộ điều khiển phần cứng ngoài thực địa, mà tập trung vào lớp quyết định thuật toán: robot phải có hành vi nhất quán khi gặp các ràng buộc về năng lượng, vật cản động và điều kiện quay về.

\section{Ghi log, thống kê và trực quan hóa}

Để phục vụ đánh giá thực nghiệm, hệ thống không chỉ hiển thị quá trình mô phỏng mà còn ghi nhận dữ liệu chạy thử. Mỗi lượt chạy có thể tạo ra bản ghi log, ảnh trạng thái cuối, dữ liệu thống kê và bản sao đầu vào. Việc lưu lại các thông tin này giúp kết quả có thể được kiểm tra lại, so sánh giữa các cấu hình và tổng hợp thành benchmark.

Các chỉ số chính được hệ thống ghi nhận được trình bày trong Bảng~\ref{tab:logged_metrics}.

\begin{table}[!htbp]
\centering
\caption{Một số chỉ số được ghi nhận sau mỗi lượt mô phỏng}
\label{tab:logged_metrics}
\small
\setlength{\tabcolsep}{4pt}
\renewcommand{\arraystretch}{1.05}
\begin{tabular}{|L{0.28\textwidth}|L{0.62\textwidth}|}
\hline
\textbf{Chỉ số} & \textbf{Ý nghĩa} \\
\hline
Tỷ lệ bao phủ & Phần trăm số ô hợp lệ đã được robot bao phủ. \\
\hline
Tổng số bước & Số bước robot đã di chuyển trên bản đồ. \\
\hline
Năng lượng đã dùng & Tổng năng lượng tiêu thụ trong quá trình mô phỏng. \\
\hline
Năng lượng di chuyển & Phần năng lượng dùng cho các bước đi sang ô kề cạnh. \\
\hline
Năng lượng quay & Phần năng lượng dùng cho thao tác đổi hướng. \\
\hline
Năng lượng còn lại & Năng lượng của robot tại thời điểm kết thúc mô phỏng. \\
\hline
Số lần quay về căn cứ & Số lần robot chuyển sang nhiệm vụ quay về căn cứ. \\
\hline
Số lần sạc & Số lần robot sạc lại tại căn cứ. \\
\hline
Kết quả nhiệm vụ & Trạng thái cuối như thành công, hoàn thành một phần hoặc thất bại. \\
\hline
Trạng thái kết thúc tại căn cứ & Cho biết robot có kết thúc ở căn cứ hay không. \\
\hline
\end{tabular}
\end{table}

Bên cạnh dữ liệu số, hệ thống sử dụng OpenCV để trực quan hóa bản đồ và quá trình di chuyển. Các thành phần như ô vật cản, ô đã bao phủ, robot, vật cản động, trạng thái năng lượng và trạng thái nhiệm vụ được hiển thị để người dùng quan sát diễn biến mô phỏng. Tuy nhiên, đánh giá chính của đồ án không chỉ dựa trên quan sát trực quan, mà dựa trên các chỉ số có thể tổng hợp và so sánh.

\section{Kết chương}

Chương này đã trình bày phương pháp đề xuất và thiết kế hệ thống mô phỏng robot bao phủ bản đồ lưới có trọng số. Hệ thống sử dụng bản đồ đầu vào gồm chi phí địa hình, vật cản tĩnh, vật cản động và căn cứ. Robot được mô hình hóa bằng trạng thái gồm vị trí, hướng, năng lượng, đường đi hiện tại và trạng thái nhiệm vụ. Bộ lập kế hoạch sử dụng Dijkstra có hướng để tìm đường theo chi phí kết hợp giữa năng lượng di chuyển và năng lượng quay.

Chương này cũng đã mô tả chiến lược chọn mục tiêu bao phủ có xét điều kiện quay về căn cứ, chính sách dự phòng quay về có điều kiện, cơ chế sạc, xử lý vật cản động, vòng lặp mô phỏng, các trạng thái nhiệm vụ và hệ thống ghi log, thống kê. Các thành phần này tạo thành một hệ thống mô phỏng hoàn chỉnh, có khả năng chạy thử, quan sát và thu thập dữ liệu để đánh giá. Trên cơ sở đó, Chương 4 sẽ phân tích sâu hơn các đặc điểm lý thuyết của phương pháp, bao gồm tác động của chi phí quay, điều kiện an toàn năng lượng, hành vi trong môi trường động và độ phức tạp của các bước lập kế hoạch chính.

\end{document}
